Vision-Language Models (VLMs) are increasingly applied to medical imaging, with
specialised fine-tunes such as MedGemma~\citep{medgemma2025} achieving strong
performance on radiology VQA tasks. Yet the \emph{mechanism} by which fine-tuning
changes these models remains opaque: does medical adaptation cause the model to
attend more carefully to the image, to rely more on language priors, or to combine
the two modalities at different decoder layers?

Many XAI methods for VLMs focus on \emph{spatial} attribution (which pixels or
image regions influence the prediction \citep{selvaraju2017gradcam,chefer2021transformer}), and do not localise, by transformer depth, where visual information becomes
causally relevant. Attention knockout answers that question: \citet{geva2023dissecting}
introduced it for text LLMs, and \citet{zhang2025crossmodal} extended per-layer
cross-modal information-flow probing to multimodal LLMs. We build directly on that line
of work and do not claim the probing mechanism as new.

Our contribution is to turn this probing into a \emph{controlled, pre-specified
before/after measurement instrument} and apply it to a matched base-vs-medical pair
(Gemma-3-4B vs.\ MedGemma-4B), asking a question prior work does not: what does medical
fine-tuning change in the image pathway, and can that change be attributed to the
encoder or the decoder? We implement the probe with lightweight PyTorch hooks.

\paragraph{Contributions.}
We contribute a pre-specified before/after instrument that applies per-layer
attention-knockout probing to Gemma-3-4B vs.\ MedGemma-4B on VQA-RAD, with a held-out
confirmation split. The probing itself follows prior work, and the instrument is
implemented with PyTorch forward hooks and no architectural changes. The open question
is whether the change lies in the encoder or the decoder: we document a large encoder
change, but leave causal attribution to a planned follow-up study.
